\section{Related Work}
\label{sec:related-work}

\paragraph{Vision-Language Pretraining.}
Recent years have witnessed the rapid development of vision-language pretraining. 
Early approaches~\cite{vlbert,vilbert,visualbert,uniter,interbert,oscar,vinvl,unimo,vlt5,m6} relied heavily on region features extracted by object detectors, which is resource\&time-consuming. 
With the increasing popularity of Vision Transformer~\cite{vit}, numerous works use Transformer to jointly learn vision-language data and demonstrate superior performance in downstream tasks~\cite{vilt,albef,meter,vlmo,fiber,lemon,blip,blip2,mplug2}. 
To facilitate alignment between vision and language modalities, researchers propose various efficient pretraining tasks.
Among them, contrastive learning is one of the most representative methods that has been widely adopted in a lot of works~\cite{clip,align,lit,florence,simclr,moco,albef,chinese_clip}.
% CLIP~\cite{clip}, the most representative foundation model based on cross-modal contrastive learning, has been widely adopted to downstream applications, e.g., image generation~\cite{dalle2,stable_diffusion}, multimodal understanding~\cite{blip,kosmos}, etc.
% The model CLIP~\cite{clip}, trained through contrastive learning, has been successfully applied in many works. For example, CLIP can be used as a foundation model to improve the performance of multimodal large language models (MLLM) ~\cite{blip2,kosmos}, and learning the features extracted by CLIP can improve the model's performance in vision tasks~\cite{fd-clip,dbot,milan,eva}. 
There also emerge some works explore the unified frameworks to handle vision-language tasks~\cite{ofa,simvlm,git,uni-perceiver,flava,coca,beit3,pali,metalm,unified_io}. \cite{ofa} adopts an encoder-decoder model to transform all vision-language tasks into generation tasks. 
\cite{coca} uses contrastive learning and text generation as the pretraining objectives, thus can be applied to image-text retrieval and vision-language generation tasks. 
\cite{beit3} employs the Multiway Transformer to process vision-language data, and discretizes images into image tokens through CLIP~\cite{clip} for joint learning with text tokens.

\paragraph{Audio-Language Pretraining.}

There is currently a significant amount of research being conducted in audio-language pretraining.
One category of these works focuses on speech-text joint pretraining. For instance, some studies propose to train a unified encoder for speech and text, which utilizes a large amount of unlabeled speech and text with masked prediction tasks and paired speech-text data to learn alignment~\cite{slam,mslam,maestro,speechlm}. 
There are also some works proposed to jointly pretrain speech and text under the encoder-decoder framework~\cite{speecht5,stpt,speechut,mmspeech}, which can be well applied to generation tasks, such as speech recognition and synthesis. 
Another category introduces cross-modal contrastive learning to audio-language pretraining~\cite{audioclip,clap,wav2clip,laion_clap}. \cite{clap} uses CNN14~\cite{panns} and BERT~\cite{bert} to extract audio and text information respectively, and conducts contrastive learning on environmental sound data. 
\cite{laion_clap} further introduces more environmental sound data and trained with HTSAT~\cite{htsat} and RoBERTa~\cite{roberta}. 
It achieves state-of-the-art results in audio downstream tasks such as audio classification and audio-text retrieval.

% There are a lot of works working on audio-language pretraining. One category of these works focuses on pretraining a unified encoder for speech and text~\cite{slam,mslam,maestro,speechlm}, which utilize a large amount of unlabeled speech and text with masked prediction tasks and paired speech-text data to learn alignment. There are also some works proposed to jointly pretrain speech and text under the encoder-decoder framework~\cite{speecht5,stpt,speechut,mmspeech}, which can be well applied to generation tasks, such as speech recognition and synthesis. In addition, some works introduce cross-modal contrastive learning into audio-language pretraining~\cite{audioclip,clap,wav2clip,laion_clap}. \cite{clap} uses CNN14~\cite{panns} and BERT~\cite{bert} to extract audio and text information respectively, and then conducts contrastive learning on environmental sound data. 
% \cite{laion_clap} further introduces more environmental sound data and trained with HTSAT~\cite{htsat} and RoBERTa~\cite{roberta}. 
% It achieves state-of-the-art results in downstream tasks such as audio classification and audio-text retrieval.

\paragraph{Vision-Audio-Language Pretraining.}
Recently researchers have begun exploring joint learning of vision, audio, and language modalities. 
\cite{vatt} employs a unified Transformer model to jointly learn video, text, and audio through cross-modal contrastive learning. 
\cite{icode,vatlm} utilize external models (e.g., VQ-VAE~\cite{vqvae} and AV-HuBERT~\cite{av_hubert}) to discretize the video and audio data, and train the models with masked prediction objectives.
~\cite{data2vec} proposes a general self-supervised learning method that does not rely on external models.
It successfully applies the method to vision, language, and audio modalities, but has not extended to multi-modal data.

% Compared to previous works, \onepeace has a flexible architecture that supports multiple modalities and general pretraining tasks that are universally applicable without external models, making it easily adaptable to diverse modalities.

% \cite{icode} uses pretrained models as modality-specific feature extractors for video, text, and audio. It discretizes video and speech data using VQ-VAE~\cite{vqvae} and Wav2Vec 2.0~\cite{wav2vec2} respectively and then employs the masked prediction task for joint learning of all modalities.
% \cite{vatlm} enhances audio representation through joint learning of video, text, and audio. It also utilizes pretrained models like AV-HuBERT~\cite{av_hubert} for data discretization and employs the masked prediction task for pretraining.
% Compared to previous works, \onepeace has strong flexibility and does not rely on extra models such as CLIP or VQ-VAE. 
Compared to previous works, \onepeace has a flexible architecture that is compatible with multiple modalities. Furthermore, the pretraining tasks of \onepeace are universally applicable without external models
Therefore, \onepeace can be easily extended to various modalities.
% Therefore, \onepeace has the potential to expand to infinite modalities.